% FORMALIZACION_REUNIDA: alineamiento dodecafasico y lector angular.
% CERTIFICADO_NUEVO: composicion exacta conservada en el suplemento local.
% La generacion previa del sello y la seleccion sectorial mantienen sus pruebas.
\section{Realización equivariante del lector angular en la incidencia de Witt}
\label{sec:ii-enlace-angular-witt}

La coordenatización angular y la incidencia excepcional conservan información diferente
del mismo desarrollo APP--TRIT--TPK. La primera expresa coordenadas de
rotación; la segunda conserva relaciones entre posiciones, soportes y
orientaciones. Su comparación requiere un transporte que preserve estas
relaciones y permita recuperar las coordenadas completas. Construimos aquí
ese transporte y demostramos su compatibilidad con las permutaciones
orientadas del triplete regional.

Los antecedentes de la composición son el sello dodecafásico \(K\), su
cara marcada \(B_K\), el sector del proyector isótipo tridimensional
\(P_3^{\mathrm{ico}}\), las ramas
compatibles de la filtración nonádica y el sistema sectorial de la
Definición~\ref{def:ii09-sistema-sectorial}. Cada uno interviene con su
estructura, y no únicamente mediante su valor escalar. Los espacios
vectoriales reales de este apartado son realizaciones posteriores de esos
sectores ya expresados. La generación del sello, el transporte de sus
coordenadas y la evaluación del lector corresponden a etapas distintas de
la composición causal.

\subsection{Triplete regional y línea de orientación}

Sea \(V\) el espacio de conteo de las tres regiones positivas de
\(\pi_{\rm HMT}\), con base ortonormal ordenada
\[
 (v_1,v_2,v_3)=(933555,393555,339555).
\]
La notación identifica cada vector de base con la región que representa.
En la ventana de prolongación de índice siete, las once ramas admitidas
se distribuyen en órbitas de tamaños \(3+3+1+1+3\) bajo la permutación
simultánea de las tres posiciones finales de las dos primeras palabras.
El sector que conserva continuación a horizonte dos consta de
\[
\begin{array}{c|c|c}
\text{rama}&\text{triple de palabras}&\text{región positiva}\\ \hline
0&200112\mid212221\mid110110&339555\\
1&200121\mid212212\mid110110&393555\\
2&200211\mid212122\mid110110&933555.
\end{array}
\]
La rama central es la que prolonga además a horizonte tres. Estos
horizontes pertenecen a la misma filtración de prolongaciones compatibles;
las cardinalidades sucesivas del soporte son \(11\), \(3\) y \(1\).

La aplicación de la cola de la primera palabra a la región positiva
sustituye coordenadamente \(1\mapsto3\) y \(2\mapsto9\), y añade
la cola regional \(555\). Por tanto, conserva la posición del símbolo
excepcional y conmuta con toda permutación de las tres posiciones. El orden
de las ramas es \((339555,393555,933555)\); el cambio desde el orden de
\(V\) es la permutación \(R=(1\,3)\).

La línea de orientación de \(V\) es \(o(V)=\det V\). La aplicación
\begin{equation}
 \Psi:\Lambda^2V\longrightarrow\operatorname{Hom}(V,\det V),
 \qquad \Psi(u\wedge v)(w)=u\wedge v\wedge w
 \label{eq:ii-witt-hodge-natural}
\end{equation}
se identifica mediante la métrica de conteo con una isometría
\(\Lambda^2V\to o(V)\otimes V\). Para
\(\omega=v_1\wedge v_2\wedge v_3\), sus valores sobre la base exterior
son
\[
 v_1\wedge v_2\mapsto\omega\otimes v_3,
 \qquad v_2\wedge v_3\mapsto\omega\otimes v_1,
 \qquad v_1\wedge v_3\mapsto-\omega\otimes v_2.
\]
En el orden angular \((12,23,13)\), el transporte conjunto de base y
orientación tiene matriz
\begin{equation}
 D=\det(R)P_R
 \begin{pmatrix}0&1&0\\0&0&-1\\1&0&0\end{pmatrix}
 =\begin{pmatrix}-1&0&0\\0&0&1\\0&-1&0\end{pmatrix}.
 \label{eq:ii-witt-orientation-D}
\end{equation}
Así,
\(D(\theta_{12},\theta_{23},\theta_{13})^{\mathsf T}
=(-\theta_{12},\theta_{13},-\theta_{23})^{\mathsf T}\).
En todo el apartado, \(P_p e_j=e_{p(j)}\) define la convención de
matrices de permutación, y el asterisco designa el adjunto de las métricas
de conteo declaradas.

\begin{proposition}[Naturalidad de la identificación orientada]
\label{prop:ii-witt-D-natural}
Para toda \(p\in S_3\),
\[
 D\Lambda^2P_p=\det(p)P_{RpR^{-1}}D,
 \qquad D^*D=I_3.
\]
\end{proposition}
\begin{proof}
La identidad
\((P_pu)\wedge(P_pv)\wedge(P_pw)=\det(p)(u\wedge v\wedge w)\)
prueba la naturalidad de \eqref{eq:ii-witt-hodge-natural}. El cambio
simultáneo del orden regional introduce la conjugación por \(R\).
Finalmente, las columnas de \(D\) son vectores coordenados con signo,
de donde \(D^*D=I_3\).
\end{proof}

La línea de orientación explica la diferencia entre el carácter
\((3,1,0)\) del triplete de posiciones y el carácter \((3,-1,0)\)
del triplete orientado, en las clases identidad, transposición y ciclo de
orden tres. Las tres transposiciones proporcionan un control explícito:
al omitir el factor determinante, la identidad anterior cambia de signo.
Un cambio global del ancla orientada se transporta, en cambio, como cambio
conjunto de coordenadas.

\subsection{Sector dodecafásico y selección variacional de la orientación}

La coordenatización del sello previamente expresado es
\begin{equation}
 K=(234,543,140,729,659,824,621,58,914,794,146,601).
 \label{eq:ii-witt-K-marked}
\end{equation}
Para explicitar su realización de incidencia, escribimos la matriz ternaria
de Paley en el calibre utilizado:
\[
 A_W=\begin{pmatrix}
 0&1&1&1&1&1\\
 1&0&1&2&2&1\\
 1&1&0&1&2&2\\
 1&2&1&0&1&2\\
 1&2&2&1&0&1\\
 1&1&2&2&1&0
 \end{pmatrix},
 \qquad C(a)=(a,aA_W),\quad a\in\mathbb F_3^6.
\]
El soporte se toma sobre las doce coordenadas de \(C(a)\). La
intersección marcada tiene las dos expresiones compatibles
\begin{equation}
 B_K=\operatorname{supp}C(010211)\cap\operatorname{supp}C(201101)
     =\{p:K_p\ge729\}=\{4,6,9,10\}.
 \label{eq:ii-witt-B-marked}
\end{equation}
La igualdad se evalúa mediante las multiplicaciones ternarias de la matriz
exhibida. En particular, la intersección se obtiene como conjunto de
posiciones, antes de utilizar su cardinal cuatro.

Para fijar completamente la representación de las doce posiciones,
consideramos \(A_5\) actuando por multiplicación izquierda sobre las
clases \(r_pC_5\), donde \(C_5=\langle(1\,2\,3\,4\,5)\rangle\).
Los representantes, escritos como listas de imágenes sobre
\(\{1,\ldots,5\}\), son
\[
\begin{array}{c|c@{\qquad}c|c}
p&r_p&p&r_p\\ \hline
1&(4,3,5,2,1)&7&(5,4,3,2,1)\\
2&(4,2,1,3,5)&8&(1,3,4,2,5)\\
3&(1,2,4,5,3)&9&(4,2,3,5,1)\\
4&(2,5,3,4,1)&10&(3,2,5,4,1)\\
5&(4,3,1,5,2)&11&(4,5,2,3,1)\\
6&(5,1,2,3,4)&12&(5,3,2,4,1).
\end{array}
\]
Denotamos esta representación de permutación por \(\rho\), y ponemos
\(V_{11}=\mathbf1^\perp\subset\mathbb R^{12}\). El proyector
isótipo tridimensional icosaédrico se determina por
\begin{equation}
 P_3^{\mathrm{ico}}=\frac3{60}\sum_{a\in A_5}\chi_3(a^{-1})\rho(a),
 \label{eq:ii-witt-projector-character}
\end{equation}
donde los valores de \(\chi_3\) sobre las clases de tipos
\(1,2^2,3,5_a,5_b\) son
\[
 \left(3,-1,0,\frac{1+\sqrt5}{2},\frac{1-\sqrt5}{2}\right).
\]
La clase \(5_a\) contiene \((1\,2\,3\,4\,5)\), y la clase
\(5_b\) contiene su cuadrado. Esta convención distingue los dos
sumandos tridimensionales. El superíndice \(\mathrm{ico}\) identifica
esta componente de la representación de \(A_5\); el proyector cíclico
\(P_3\) que interviene en~\eqref{eq:det} actúa sobre el transporte de doce
posiciones con su definición propia. La evaluación finita de
\eqref{eq:ii-witt-projector-character} satisface
\(P_3^{\mathrm{ico}}=(P_3^{\mathrm{ico}})^*=(P_3^{\mathrm{ico}})^2\), \(P_3^{\mathrm{ico}}\mathbf1=0\) y
\(\operatorname{Tr}P_3^{\mathrm{ico}}=3\).

Los dos vectores marcados del sector son
\[
 b=P_3^{\mathrm{ico}}\left(\mathbf1_{B_K}-\frac13\mathbf1\right),
 \qquad u=P_3^{\mathrm{ico}}(K-\overline K\mathbf1),
 \qquad \overline K=\frac1{12}\sum_{p=1}^{12}K_p,
 \qquad \|b\|^2=1.
\]
Para cada uno de los diez subgrupos de orden tres \(C_3\subset A_5\),
sea \(H=N_{A_5}(C_3)\cong S_3\). El carácter
\(\operatorname{sgn}_H\) toma valor \(+1\) sobre \(C_3\) y
\(-1\) sobre \(H\setminus C_3\). Es el carácter del cociente
\(H/C_3\), mientras el signo ambiental de los elementos de \(A_5\)
es siempre positivo. Se define
\begin{equation}
 e_{{\rm sgn},H}=\frac16\sum_{a\in H}\operatorname{sgn}_H(a)\rho(a),
 \qquad \mathcal E_B(H)=\|e_{{\rm sgn},H}b\|^2.
 \label{eq:ii-witt-energy}
\end{equation}

\begin{proposition}[Normalizador y orientación determinados por el sello]
\label{prop:ii-witt-variational-selection}
La regla de máximo de \eqref{eq:ii-witt-energy} selecciona un único
normalizador \(H_*\). Su energía y la segunda energía distinta son
\[
 \frac23+\frac{2\sqrt5}{15},\qquad
 \frac13+\frac{2\sqrt5}{15},
\]
respectivamente; la separación es \(1/3\). La proyección de \(u\)
selecciona el mismo normalizador. Dentro de \(H_*\), la maximización de
\(s(a)=\langle b,\rho(a)u\rangle\), por separado entre los elementos
de orden tres y las involuciones, determina
\[
 c_*=(1,2,4,5,3),\qquad r_*=(2,1,5,4,3)
\]
en notación de imágenes sobre cinco puntos. Se cumple
\(r_*c_*r_*=c_*^{-1}\).
\end{proposition}
\begin{proof}
La evaluación de las sumas finitas anteriores en \(\mathbb Q(\sqrt5)\)
da la siguiente distribución de las diez energías:
\[
\begin{array}{c|c@{\qquad}c|c}
\text{energía}&\text{multiplicidad}&\text{energía}&\text{multiplicidad}\\ \hline
2/3+2\sqrt5/15&1&2/3-2\sqrt5/15&1\\
1/3+2\sqrt5/15&2&1/3-2\sqrt5/15&2\\
1/6+\sqrt5/30&2&1/6-\sqrt5/30&2.
\end{array}
\]
La mayor pertenece al subgrupo generado por \(c_*\). La siguiente se
obtiene restando \(1/3\). La misma evaluación con \(u\) tiene máximo
único en ese subgrupo, de valor
\(2526869/12+2813609\sqrt5/30\).
La puntuación de las dos orientaciones del ciclo es
\[
 s(c_*)=\frac{1341}{4}+\frac{2409\sqrt5}{20},\qquad
 s(c_*^{-1})=\frac{357}{2}+\frac{1853\sqrt5}{10}.
\]
Su diferencia es \(627/4-1297\sqrt5/20>0\), como se comprueba
comparando los cuadrados positivos. Las puntuaciones de las tres
involuciones son
\[
\begin{array}{c|c}
\text{lista de imágenes}&s(a)\\ \hline
(2,1,3,5,4)&-1341/4-2409\sqrt5/20\\
(2,1,4,3,5)&-1659/4-2667\sqrt5/20\\
(2,1,5,4,3)&-357/2-1853\sqrt5/10.
\end{array}
\]
El máximo único corresponde a \(r_*\). La composición de las dos listas de
imágenes verifica la relación diédrica. Todas estas operaciones usan
\(K\), \(B_K\), \(P_3^{\mathrm{ico}}\) y las acciones finitas ya especificadas.
\end{proof}

Escribimos \(\vartheta:S_3\to H_*\) para el isomorfismo determinado
por \(\vartheta((1\,2\,3))=c_*\) y
\(\vartheta((1\,3))=r_*\). La reflexión de posiciones fija la rama
central que alcanza horizonte tres. Esta orientación completa la
identificación regional antes de cualquier evaluación de ángulos.

\subsection{Entrelazador de Reynolds y normalización polar}

En la base de ramas, sea \(\rho_o(g)=\det(g)P_g\) y
\(w=(0,1,0)^{\mathsf T}\). El operador de Reynolds definido por suma
sobre el grupo es
\begin{equation}
 A_R=\sum_{g\in S_3}\rho_o(g)\,|w\rangle\langle b|\,
                       \rho(\vartheta(g))^{-1}.
 \label{eq:ii-witt-reynolds}
\end{equation}
El símbolo \(A_R\) distingue este operador del ángulo
\(A_{\rm deg}\). Cada sumando tiene escala unitaria; la normalización
métrica se efectúa mediante su factor polar.

\begin{proposition}[Isometría polar del sector orientado]
\label{prop:ii-witt-polar}
Se cumplen
\[
 A_R=A_RP_3^{\mathrm{ico}},\qquad
 \rho_o(g)A_R=A_R\rho(\vartheta(g)),\qquad
 G=A_RA_R^*>0,\qquad A_R^*G^{-1}A_R=P_3^{\mathrm{ico}}.
\]
Con \(E=\mathbf1\mathbf1^{\mathsf T}/3\),
\begin{equation}
 G=\lambda_sE+\lambda_t(I_3-E),\qquad
 \lambda_s=8+\frac{8\sqrt5}{5},\quad
 \lambda_t=\frac12-\frac{\sqrt5}{5}>0.
 \label{eq:ii-witt-gram-spectrum}
\end{equation}
En consecuencia,
\[
 U_{GK}=G^{-1/2}A_R:P_3^{\mathrm{ico}}V_{11}\longrightarrow o(V)\otimes V
\]
es una isometría sobreyectiva:
\(U_{GK}^*U_{GK}=P_3^{\mathrm{ico}}\) y \(U_{GK}U_{GK}^*=I_3\).
\end{proposition}
\begin{proof}
El proyector \(P_3^{\mathrm{ico}}\) conmuta con \(\rho\), y \(b=P_3^{\mathrm{ico}}b\);
por ello \(A_R=A_RP_3^{\mathrm{ico}}\). Reindexar la suma de Reynolds prueba la
equivarianza. La multiplicación de la suma finita da entradas diagonales
\(3+2\sqrt5/5\) y entradas no diagonales
\(5/2+3\sqrt5/5\) en \(G\). Las proyecciones \(E\) e \(I_3-E\)
proporcionan \eqref{eq:ii-witt-gram-spectrum}. Ambos autovalores son
positivos; para el menor basta \(25>20\). Por tanto \(A_R\) tiene
rango tres. Su espacio fila está contenido en el sector de rango tres
de \(P_3^{\mathrm{ico}}\), y coincide con él. El proyector ortogonal sobre ese espacio
fila es \(A_R^*(A_RA_R^*)^{-1}A_R\), lo que da la última identidad.
Las dos ecuaciones del factor polar se siguen por sustitución.
\end{proof}

La raíz positiva conserva la equivarianza porque ambos proyectores
espectrales conmutan con la representación orientada. La unicidad del
factor polar corresponde al operador de Reynolds y a la orientación
determinados por las reglas precedentes.

\subsection{Incidencia punto--hexada y cambio conjunto de calibre}

Sea \(\mathcal W\) el diseño \(S(5,6,12)\) de 132 hexadas obtenido
de los soportes de peso seis del código ternario \(C(a)\). La matriz
de incidencia es
\[
 I_{H,p}=\mathbf1_{p\in H},\qquad
 (Iz)_H=\sum_{p\in H}z_p.
\]
Los bloques se reconstruyen enumerando \(a\in\mathbb F_3^6\),
reteniendo peso seis y reuniendo soportes iguales. La verificación
finita de que cada subconjunto de cinco puntos pertenece a una única
hexada fija la propiedad de diseño utilizada a continuación.

\begin{lemma}[Identidades métrica y espectral de la incidencia]
\label{lem:ii-witt-incidence-identities}
Sean \(J_{12}\) y \(J_{132,12}\) las matrices de unos indicadas por
sus dimensiones, y sea
\[
 (G_{\mathcal W})_{H,H'}=\binom{|H\cap H'|}{4}.
\]
Entonces
\begin{equation}
 I^*I=36I_{12}+30J_{12},\qquad
 G_{\mathcal W}I=30I+15J_{132,12}.
 \label{eq:ii-witt-incidence-identities}
\end{equation}
En particular, \(U_W=I/6\) es una isometría sobre
\(V_{11}=\mathbf1^\perp\), con imagen contenida en
\(E_{30}=\ker(G_{\mathcal W}-30I_{132})\).
\end{lemma}
\begin{proof}
La propiedad de diseño da 66 hexadas por punto y 30 por par de puntos:
\[
 \frac{\binom{11}{4}}{\binom54}=66,
 \qquad \frac{\binom{10}{3}}{\binom43}=30.
\]
Éstas son las entradas diagonales y no diagonales de \(I^*I\).
Para la segunda identidad se cuentan los pares \((T,H')\), donde
\(T\subset H\cap H'\), \(|T|=4\) y \(p\in H'\).
Si \(p\in H\), diez tetradas \(T\subset H\) contienen \(p\),
y cada una pertenece a cuatro hexadas; las otras cinco tetradas,
unidas a \(p\), determinan una hexada cada una. El resultado es
\(10\cdot4+5=45\). Si \(p\notin H\), cada una de las quince
tetradas determina con \(p\) una hexada, y el resultado es 15.
Así se obtiene \(30I_{H,p}+15\). Las matrices de unos se anulan sobre
\(V_{11}\), lo que prueba la isometría y la pertenencia espectral.
\end{proof}

El cambio de calibre de las doce posiciones es
\begin{equation}
 \sigma=(1,5,2,6,7,4,3,10,11,9,8,12),
 \label{eq:ii-witt-sigma}
\end{equation}
escrito como lista de imágenes con índices desde uno. Si
\(\rho_{\rm perm}\) es la acción puntual subyacente a \(\rho\), las
seis permutaciones
\[
 \tau_g=\sigma\rho_{\rm perm}(\vartheta(g))\sigma^{-1},\qquad g\in S_3,
\]
preservan \(\mathcal W\). Sus dos generadores tienen listas de imágenes
\[
\begin{aligned}
 \tau_{(123)}&=(3,6,5,2,1,4,9,12,11,8,7,10),\\
 \tau_{(13)}&=(2,1,4,3,6,5,8,7,10,9,12,11).
\end{aligned}
\]
La afirmación se comprueba sobre el diseño finito: se aplican estas dos
permutaciones generadoras a los soportes
\(C(a)\) de peso seis y se verifica que sus imágenes recorren los
mismos 132 soportes. Los productos de estos generadores dan las otras
cuatro permutaciones. La acción inducida \(Q_{\tau_g}\) sobre hexadas
satisface
\begin{equation}
 I P_{\tau_g}=Q_{\tau_g}I,
 \qquad
 P_\sigma\rho(\vartheta(g))=P_{\tau_g}P_\sigma.
 \label{eq:ii-witt-recalibration}
\end{equation}
La primera identidad expresa exactamente la equivalencia
\(p\in H\Longleftrightarrow\tau_g(p)\in\tau_g(H)\).

Este cambio transporta conjuntamente posiciones, código, incidencia,
origen, cámara y orientación del sello. En el calibre inicial, cinco
elementos de este subgrupo de orden seis no preservan el diseño fijo.
La conjugación \eqref{eq:ii-witt-sigma} proporciona el calibre compatible
y mantiene ese control sobre el etiquetado inicial. El procedimiento
exhibe un cambio de calibre; su unicidad no interviene en la prueba.

\subsection{Composición isométrica y recuperación de coordenadas}

\begin{theorem}[Transporte angular dodecafásico--Witt]
\label{thm:ii-witt-angular-transport}
La composición
\begin{equation}
 \boxed{F=\frac16 I P_\sigma A_R^*G^{-1/2}D:
        \mathbb R^3_{\rm ang}\longrightarrow E_{30}}
 \label{eq:ii-witt-F}
\end{equation}
es una isometría. Su acción coordenada y su equivarianza son
\begin{align}
 (F\boldsymbol\theta)_H
 &=\frac16\sum_{p\in H}
       (P_\sigma A_R^*G^{-1/2}D\boldsymbol\theta)_p,
       \label{eq:ii-witt-F-coordinates}\\
 F X_g&=Q_{\tau_g}F,
 \qquad X_g=\Lambda^2P_{R^{-1}gR}.
       \label{eq:ii-witt-F-equivariance}
\end{align}
\end{theorem}
\begin{proof}
Las columnas de \(A_R^*\) pertenecen a \(P_3^{\mathrm{ico}}V_{11}\), y
\(P_\sigma\) conserva \(V_{11}\). El
Lema~\ref{lem:ii-witt-incidence-identities} demuestra por tanto la
pertenencia de la imagen de \(F\) a \(E_{30}\). Además,
\begin{align*}
 F^*F
 &=D^*G^{-1/2}A_RP_\sigma^*
       \frac{I^*I}{36}P_\sigma A_R^*G^{-1/2}D\\
 &=D^*G^{-1/2}A_RA_R^*G^{-1/2}D
  =D^*D=I_3.
\end{align*}
La fórmula coordenada desarrolla la multiplicación por la incidencia.
Para la equivarianza se componen, en su orden, la naturalidad de \(D\),
la equivarianza de \(A_R^*\), la conmutación de \(G^{-1/2}\) con
\(\rho_o\) y las dos identidades de
\eqref{eq:ii-witt-recalibration}. Esto produce
\eqref{eq:ii-witt-F-equivariance} para cada \(g\in S_3\).
\end{proof}

\begin{proposition}[Covariancia del calibre completo]
\label{prop:ii-witt-gauge-covariance}
Sean \(S,T,V_o\) cambios ortogonales de coordenadas en el espacio
dodecafásico, el espacio puntual de Witt y el sector orientado,
respectivamente. Sea \(Q_T\) el cambio correspondiente de coordenadas
en el espacio de incidencia; para un reetiquetado del diseño es la
permutación inducida de hexadas. Si se transportan conjuntamente
\[
\begin{gathered}
 A_R'=V_oA_RS^*,\quad G'=V_oGV_o^*,\quad
 P_\sigma'=TP_\sigma S^*,\\
 D'=V_oD,\qquad U_W'=Q_TU_WT^*,
\end{gathered}
\]
la composición satisface \(F'=Q_TF\).
\end{proposition}
\begin{proof}
La unicidad de la raíz positiva da
\((G')^{-1/2}=V_oG^{-1/2}V_o^*\). Sustituyendo las cinco
transformaciones en
\(F'=U_W'P_\sigma'(A_R')^*(G')^{-1/2}D'\), los factores
\(T^*T\), \(S^*S\) y \(V_o^*V_o\) se cancelan por pares, y queda
\(Q_TU_WP_\sigma A_R^*G^{-1/2}D=Q_TF\).
\end{proof}

La covariancia precedente describe un cambio de coordenadas de todos los
datos marcados. Su aplicación a una transformación cardinal activa requiere
la identidad correspondiente entre las acciones sobre el estado; ambas
nociones conservan así sus dominios propios.

\subsection{Realización fiel del sistema sectorial angular}

Se utiliza la coordenatización \(x_\theta\) de
\eqref{eq:ii-angular-chart},
\[
 x_\theta=\begin{pmatrix}A_{\rm deg}\\C^*_{\rm deg}/6\\
                              (180/\pi_{\rm HMT})\Delta_4\end{pmatrix}.
\]
La componente \(h_\theta=C^*_{\rm deg}/6\) distribuye el componente
impar del registro entre seis parejas antipodales. La
Definición~\ref{def:ii09-sistema-sectorial} fija las tres reglas
\[
 R_{12}=(2,-p,bq),\quad R_{23}=(0,q,-p^2/2),\quad
 R_{13}=(1,-pb,-(o_8-h_6)/t)
\]
para \((b,p,q,h_6,o_8,t)=(4,5,7,6,8,3)\). La cara de cardinal
cuatro, el orden pentádico, la coordenada del pivote, la diferencia
hexádico--octádica y el alfabeto trítico mantienen tipos distintos.
La normalización orbital bidireccional de
\ref{prop:ii09-semicuadrado-orbital} conserva el término \(p^2/2\),
incluida la contribución de los pares fijos.

Su evaluación, ya demostrada en el
Teorema~\ref{thm:ii09-matriz-sectorial}, da
\[
 \mathsf M_{\rm CKM}=
 \begin{pmatrix}2&-5&28\\0&7&-25/2\\1&-20&-2/3\end{pmatrix},
 \qquad\det\mathsf M_{\rm CKM}=-\frac{3857}{6}.
\]
El transporte \(F\) actúa sobre las coordenadas expresadas por esas
reglas; la elección anterior de los funcionales conserva su lugar en la
construcción sectorial.

\begin{corollary}[Lector angular en el módulo de incidencia]
\label{cor:ii-witt-composed-reader}
La aplicación
\begin{equation}
 \mathcal C=F\mathsf M_{\rm CKM}
 \label{eq:ii-witt-composed-reader}
\end{equation}
satisface
\[
 \mathcal C^*\mathcal C=\mathsf M_{\rm CKM}^*\mathsf M_{\rm CKM},
 \qquad
 \mathsf M_{\rm CKM}^{-1}F^*\mathcal C=I_3.
\]
Es una isometría del espacio de \(x_\theta\) provisto de la métrica
\(\mathsf M_{\rm CKM}^*\mathsf M_{\rm CKM}\) sobre su imagen.
\end{corollary}
\begin{proof}
Se sustituye \(F^*F=I_3\) en ambos productos. El determinante no
nulo de \(\mathsf M_{\rm CKM}\) permite utilizar su inversa,
exhibida en \eqref{eq:ii-ckm-inverse}. La positividad de la métrica
se sigue de
\(x^*\mathsf M_{\rm CKM}^*\mathsf M_{\rm CKM}x
=\|\mathsf M_{\rm CKM}x\|^2>0\) para \(x\ne0\).
\end{proof}

La isometría \(F\) conserva la métrica de conteo de
\(\boldsymbol\theta=\mathsf M_{\rm CKM}x_\theta\); su transporte
a las coordenadas \(x_\theta\) da la métrica del corolario. Esta
distinción evita identificar una base de generadores angulares con una
base ortonormal de sus parámetros.

\subsection{Certificación exacta y composición unitaria finita}

Las raíces positivas del factor polar pueden comprobarse mediante sus
proyectores, conservando la aritmética exacta. Definimos
\begin{equation}
 B=\frac16IP_\sigma A_R^*D,\qquad
 H=D^*GD=B^*B,\qquad L=H^{-1}B^*.
 \label{eq:ii-witt-radical-free}
\end{equation}
Las identidades
\[
 LB=I_3,\qquad L(B\mathsf M_{\rm CKM})=\mathsf M_{\rm CKM},\qquad
 \mathsf M_{\rm CKM}^{-1}L(B\mathsf M_{\rm CKM})=I_3
\]
se evalúan en \(\mathbb Q(\sqrt5)\). Si \(E_D=D^*ED\),
\[
 H^{-1/2}=\lambda_s^{-1/2}E_D+
          \lambda_t^{-1/2}(I_3-E_D),\qquad F=BH^{-1/2}.
\]
Los productos \(E_D^2=E_D\), \(E_D(I_3-E_D)=0\) y la
descomposición de \(H\) prueban esta fórmula sin aproximar sus
radicales. Las 132 filas de \(B\) quedan determinadas por
\eqref{eq:ii-witt-F-coordinates} y por la regla de construcción de
hexadas. El certificado ejecutable de la entrega conserva esas filas,
las de \(B\mathsf M_{\rm CKM}\), la inversa izquierda \(L\),
los proyectores y las etiquetas de cada hexada. Sus 906 controles
exactos comprueban esta composición y sus identidades de orientación,
incidencia y recuperación.

La conversión a radianes se efectúa antes de la realización trigonométrica:
\begin{equation}
 \boldsymbol\vartheta=
 \mathsf M_{\rm CKM}
 \begin{pmatrix}(\pi_{\rm HMT}/180)A_{\rm deg}\\
 (\pi_{\rm HMT}/180)C^*_{\rm deg}/6\\\Delta_4\end{pmatrix},
 \qquad
 \delta=\frac{\pi_{\rm HMT}}{180}\,9A_{\rm deg}.
 \label{eq:ii-witt-angular-units}
\end{equation}
La realización finita preserva el orden
\(V_{\rm CKM}=R_{23}U_{13}(\delta)R_{12}\) de
\eqref{eq:ii-ckm-factorization}. Para precisar la información adicional
de la fase compleja, sea \(E_{ab}\) la matriz elemental y definamos
\[
 J_{13}=E_{13}-E_{31},\qquad S_{13}=E_{13}+E_{31},\qquad
 X_{13}(\delta)=\cos\delta\,J_{13}-i\sin\delta\,S_{13}.
\]
Entonces \(X_{13}(\delta)^*=-X_{13}(\delta)\) y
\(U_{13}(\delta)=\exp(\vartheta_{13}X_{13}(\delta))\).
En efecto, \(X_{13}(\delta)^2=-(E_{11}+E_{33})\), y su exponencial
reproduce el bloque unitario \(13\) de
\eqref{eq:ii-ckm-factorization}. El transporte exterior conserva los
generadores orientados; la componente imaginaria simétrica conserva
la información de fase. El producto ordenado de los tres factores
finitos mantiene estas dos contribuciones y su orden de composición.

\subsection{Resultado estructural y alcance de la composición}

La construcción proporciona un transporte isométrico equivariante desde
las tres coordenadas angulares hasta la incidencia de Witt, junto con una
reconstrucción explícita de \(A_{\rm deg}\), \(C^*_{\rm deg}/6\) y
\(\Delta_{\rm deg}\). Conserva la orientación exterior, el sector
dodecafásico, la cara marcada del sello y el cambio conjunto de calibre.
En particular, la realización de incidencia representa el mismo contenido
angular en otro soporte y permite recuperar sus tres coordenadas.

La generación anterior de \(K\), la selección de los funcionales
sectoriales y la identificación física de las amplitudes conservan sus
demostraciones propias. Las identidades de recuperación aquí probadas
establecen la fidelidad de la composición; su uso posterior mantiene
esos antecedentes explícitos. El desarrollo complementario de los planos
cuárticos de memoria permanece como otra construcción del estado, y no
se interpone como requisito adicional del alineamiento realizado en
\eqref{eq:ii-witt-F}.
